\documentclass[a4paper,11pt]{article}
\usepackage[utf8]{inputenc}
\usepackage[T1]{fontenc}
\usepackage[brazil]{babel}
% Ajuste de geometria para acomodar o cabeçalho grande
\usepackage[left=2cm, right=2cm, top=3.5cm, bottom=2.5cm, headheight=50pt]{geometry}
\usepackage{xcolor}
\usepackage{graphicx}
\usepackage{tcolorbox}
\usepackage{helvet}
\usepackage{tikz}
\usepackage{fancyhdr}
\usepackage{array}
\usepackage{colortbl}
\usepackage{hyperref}

% --- CONFIGURAÇÕES DE LAYOUT E BIBLIOTECAS ---
\tcbuselibrary{skins, raster, breakable}
\renewcommand{\familydefault}{\sfdefault}
\setlength{\parindent}{0pt}
\setlength{\parskip}{0.8em}

% --- PALETA DE CORES INFODENTAL (REBRANDING) ---
\definecolor{infoDentalTeal}{RGB}{26, 62, 71}     % Fundo Escuro (Dark Teal)
\definecolor{infoDentalGreen}{RGB}{79, 208, 104}  % Acento Verde (Bright Green)
\definecolor{infoDentalRed}{RGB}{231, 76, 60}     % Vermelho Alerta
\definecolor{infoDentalOrange}{RGB}{243, 156, 18} % Laranja (Atenção/Suave)
\definecolor{infoDentalGray}{RGB}{240, 242, 245}  % Fundo Suave das Tabelas
\definecolor{textDark}{RGB}{44, 62, 80}           % Texto Leitura

% --- CABEÇALHO E RODAPÉ PERSONALIZADOS ---
\pagestyle{fancy}
\fancyhf{} % Limpa tudo
\renewcommand{\headrulewidth}{0pt}

% Desenho do Cabeçalho (Barra Teal no Topo)
\fancyhead[C]{%
    \begin{tikzpicture}[remember picture, overlay]
        \fill[infoDentalTeal] (current page.north west) rectangle ([yshift=-3cm]current page.north east);
        \node[anchor=west, text=white, font=\bfseries\Huge] at ([xshift=2cm, yshift=-1.5cm]current page.north west) {InfoDental \textcolor{infoDentalGreen}{Intelligence}};
        \node[anchor=east, text=white, align=right, font=\small] at ([xshift=-2cm, yshift=-1.5cm]current page.north east) {
            \textbf{RELATÓRIO DE AUDITORIA FORENSE}\\
            PROTOCOLO: SHOPPER-MG-TX \\
            DATA: \today
        };
        % Linha de destaque verde
        \fill[infoDentalGreen] (current page.north west) ++(0,-2.9cm) rectangle ++(\paperwidth, 0.1cm);
    \end{tikzpicture}
}

% Desenho do Rodapé
\fancyfoot[C]{%
    \begin{tikzpicture}[remember picture, overlay]
        \fill[infoDentalGray] (current page.south west) rectangle ([yshift=1.5cm]current page.south east);
        \node[text=infoDentalTeal, font=\footnotesize] at ([yshift=0.75cm]current page.south) {Confidencial - Gerado por InfoDental AI Engine - Página \thepage};
    \end{tikzpicture}
}

% --- ÍCONES MANUAIS ---
\newcommand{\iconAlert}{%
\begin{tikzpicture}[baseline=-3pt]
    \fill[infoDentalOrange] (0,0) circle (0.25);
    \node[white, font=\bfseries\small] at (0,0) {!};
\end{tikzpicture}%
}

\newcommand{\iconCheck}{%
\begin{tikzpicture}[baseline=-3pt]
    \fill[infoDentalGreen] (0,0) circle (0.25);
    \draw[infoDentalTeal, very thick] (-0.1,-0.05) -- (0,0.1) -- (0.15,-0.15);
\end{tikzpicture}%
}

\begin{document}

% --- BADGE DE STATUS ---
\begin{flushright}
    \begin{tikzpicture}
        \node[fill=infoDentalOrange!10, draw=infoDentalOrange, rounded corners, inner sep=5pt, line width=1pt] {
            \textbf{\textcolor{infoDentalOrange}{STATUS: LATÊNCIA E FRICÇÃO DE PROCESSO}}
        };
    \end{tikzpicture}
\end{flushright}

% --- SEÇÃO 1: O CENÁRIO ---
\begin{tcolorbox}[
    enhanced,
    title=\textbf{1. O CENÁRIO DO TESTE (A ISCA)},
    colframe=infoDentalTeal, colback=white, coltitle=white,
    boxrule=0.5pt, drop shadow
]
    \textbf{Script Aplicado:} ``O Viajante com Pressa'';

    \textbf{Perfil do Avatar:} Visitante de Goiás (DDD 62) a trabalho em Teixeiras/MG;

    \textbf{Contexto:} Quebra de faceta/dente frontal (Estética = Urgência Emocional);

    \textbf{Gatilhos de Urgência:} ``Reuniões importantes'' + ``Pago particular'' + ``Só fico até terça'';

    \textbf{Telefone Contactado} +55 31 9 9540-9898


\end{tcolorbox}

% --- SEÇÃO 2: A LINHA DO TEMPO ---
\vspace{0.5cm}
\section*{\textcolor{infoDentalTeal}{2. ANÁLISE MINUTO A MINUTO}}

% Tabela customizada com tcolorbox
\begin{tcolorbox}[enhanced, frame hidden, colback=white]
    \renewcommand{\arraystretch}{1.5}
    \begin{tabular}{ | p{0.12\textwidth} | p{0.20\textwidth} | p{0.58\textwidth} | }
        \hline
        \rowcolor{infoDentalTeal}
        \textbf{\textcolor{white}{HORA}} & \textbf{\textcolor{white}{AGENTE}} & \textbf{\textcolor{white}{ANÁLISE DE FLUXO E TEMPO}} \\
        \hline
        \textbf{09:55} & Paciente (Lead) &
        \small Disparo inicial relatando dor aguda, quebra estética e urgência de tempo. \\
        \hline
        \textbf{10:18} & Clínica (Humano) &
        \textbf{Tempo de Resposta: 23 Minutos.} \newline
        \textit{``Tudo bem ?? Desculpa a demora. O celular está cheio de mensagens. Poderia me mandar uma foto por gentileza??''} \\
        \hline
    \end{tabular}
\end{tcolorbox}

% PONTO DE ATENÇÃO 1: GAP DE TEMPO
\begin{tcolorbox}[colback=infoDentalOrange!5, colframe=infoDentalOrange, boxrule=0.5pt, arc=0mm, left=2mm, right=2mm]
    \textbf{\textcolor{infoDentalOrange}{PONTO DE ATENÇÃO \#1: ``GAP DE RECEITA''}}
    \footnotesize
    O tempo de espera de \textbf{23 minutos} para uma emergência estética é crítico. Estatisticamente, após 5 minutos, a chance de o paciente contatar um concorrente sobe 400\%. A própria atendente valida o erro pedindo ``Desculpa a demora''.
\end{tcolorbox}

% PONTO DE REFLEXÃO ESTRATÉGICA (FOTO) - NOVO!
\begin{tcolorbox}[colback=infoDentalTeal!5, colframe=infoDentalTeal, boxrule=0.5pt, arc=0mm, left=2mm, right=2mm]
    \textbf{\textcolor{infoDentalTeal}{REFLEXÃO ESTRATÉGICA: A BARREIRA DA FOTO}}
    \footnotesize
    A solicitação de foto para triagem é tecnicamente válida, mas comercialmente perigosa para este perfil.
    O paciente estava \textbf{no trabalho} (ambiente corporativo). Pedir uma ``selfie intraoral'' nesse contexto gera constrangimento e fricção.
    \textit{Observe que o paciente resistiu: ``Cara, tá difícil tirar foto aqui''.}
    \textbf{O Risco:} Quantos leads tímidos ou ocupados desistem de responder quando a foto é solicitada antes do acolhimento?
\end{tcolorbox}

\begin{tcolorbox}[enhanced, frame hidden, colback=white]
    \renewcommand{\arraystretch}{1.5}
    \begin{tabular}{ | p{0.12\textwidth} | p{0.20\textwidth} | p{0.58\textwidth} | }
        \hline
        \textbf{10:21} & Clínica (Humano) &
        \textit{``Se eu conseguir um horário para você agora você tem disponibilidade de vim ,agora ?''} \\
        \hline
    \end{tabular}
\end{tcolorbox}

% PONTO DE ATENÇÃO 2: GRAMÁTICA
\begin{tcolorbox}[colback=infoDentalOrange!5, colframe=infoDentalOrange, boxrule=0.5pt, arc=0mm, left=2mm, right=2mm]
    \textbf{\textcolor{infoDentalOrange}{PONTO DE ATENÇÃO \#2: GRAMÁTICA E PONTUAÇÃO}}
    \footnotesize
    1. \textbf{Gramática:} Uso incorreto do verbo (``disponibilidade de \textbf{vim}'' em vez de ``vir'').
    2. \textbf{Pontuação:} Vírgula deslocada (``vim ,agora ?'') e duplicidade de interrogação (``Tudo bem ??'').
    Para um paciente High-Ticket, a comunicação escrita deve refletir a mesma excelência técnica do tratamento.
\end{tcolorbox}

% --- SEÇÃO 3: DIAGNÓSTICO VISUAL ---
\vspace{0.2cm}
\begin{tcolorbox}[
    enhanced,
    colback=infoDentalGray,
    colframe=white,
    title=\textcolor{infoDentalTeal}{\textbf{3. DIAGNÓSTICO OPERACIONAL}},
    coltitle=infoDentalTeal,
    fonttitle=\large\bfseries,
    frame hidden
]
    \begin{itemize}
        \item[\iconAlert] \textbf{O ``Speed to Lead'':} Embora a equipe tenha conseguido agendar, o delay inicial de 23 minutos colocou a venda em risco.
        \item[\iconAlert] \textbf{Sobrecarga Visível:} A frase ``O celular está cheio de mensagens'' revela que sua recepção está competindo entre atendimento presencial e digital.
    \end{itemize}
\end{tcolorbox}

% --- SEÇÃO 4: A SOLUÇÃO INFODENTAL ---
\vspace{0.5cm}

\begin{tcolorbox}[
    skin=enhanced,
    colback=infoDentalTeal,
    colframe=infoDentalGreen,
    boxrule=1pt,
    arc=4mm,
    shadow={2mm}{-2mm}{0mm}{black!30},
    title=\textbf{\large INFODENTAL: INTELIGÊNCIA CLÍNICA VERTICAL},
    coltitle=infoDentalGreen,
    fonttitle=\bfseries
]
    \color{white}
    Diferente de ``Chatbots'' que apenas respondem, o InfoDental é um \textbf{Sistema Vertical} que atua em todas as camadas do processo clínico:

    \vspace{0.3cm}

    \begin{itemize}
        \item[\iconCheck] \textbf{Camada de Infraestrutura (Velocidade):} Resposta instantânea (< 30 segundos) 24/7. Elimina o ``vácuo'' de 23 minutos.
        \item[\iconCheck] \textbf{Camada de Comunicação (Qualidade):} Português impecável. O sistema entende o contexto (paciente no trabalho) e foca no agendamento, deixando a foto como opcional para não perder a venda.
        \item[\iconCheck] \textbf{Camada de Processo (Alívio):} Remove a sobrecarga da equipe (``celular cheio''), permitindo foco no presencial.
    \end{itemize}

    \vspace{0.5cm}
    \centering
    \begin{tcolorbox}[colback=infoDentalGreen, colframe=infoDentalGreen, width=12cm, halign=center, arc=4mm, boxsep=15pt, top=8pt, bottom=8pt]
        \href{https://infodental.dental}{\textbf{\Large\textcolor{infoDentalTeal}{PROFISSIONALIZAR MEU SISTEMA}}}
    \end{tcolorbox}
    \vspace{0.3cm}
\end{tcolorbox}

\end{document}